\section{引言：推理的重要性}

本节课的主题是\textbf{推理}（Inference）：给定一个已经训练好的固定模型，根据输入的提示（prompt）生成响应。推理看似简单，但其效率优化涉及硬件特性、算法设计和系统工程等多个层面，是大语言模型实际应用中最关键的环节之一。

\begin{importantbox}{推理的定义}
\textbf{推理}（Inference）：给定一个\textbf{固定模型}，根据提示生成响应。与训练不同，推理不更新模型参数，但需要反复执行，因此效率至关重要。
\end{importantbox}

\subsection{推理的应用场景}

推理不仅仅是``搭建一个聊天机器人''那么简单，它在以下多个场景中扮演核心角色：

\begin{enumerate}
    \item \textbf{实际使用}：聊天机器人、代码补全（如 Cursor）、批量数据处理等
    \item \textbf{模型评估}：在指令遵循等任务上评估模型性能，本质上也是推理
    \item \textbf{测试时计算}（Test-time Compute）：让模型``思考''更多步骤再输出最终答案——思考本身就是生成 token
    \item \textbf{强化学习训练}：采样响应并根据奖励评分，采样过程即推理
\end{enumerate}

\begin{knowledgebox}{推理的规模有多大？}
一些数据帮助理解推理的规模：
\begin{itemize}
    \item OpenAI 每天生成超过 1000 亿个词
    \item Cursor 每天生成约 10 亿行被用户接受的代码
\end{itemize}
训练是一次性成本，而推理会被反复执行无数次，因此推理成本在总成本中的占比日益增长。
\end{knowledgebox}

\subsection{推理的核心度量指标}

\begin{importantbox}{推理性能的三个核心指标}
\begin{itemize}
    \item \textbf{首 token 延迟}（TTFT, Time-to-First-Token）：用户等待第一个生成 token 出现的时间。对交互式应用至关重要
    \item \textbf{延迟}（Latency, 秒/token）：每个 token 的生成速度。影响用户的流式阅读体验
    \item \textbf{吞吐量}（Throughput, tokens/秒）：系统整体每秒生成的 token 数。对批处理应用尤为重要
\end{itemize}
\end{importantbox}

\begin{warningbox}{高吞吐量不等于低延迟}
吞吐量衡量的是系统整体的处理能力，而延迟关注的是单个用户的体验。一个系统可以有很高的吞吐量（同时处理大量请求），但某些请求的延迟可能很高。这两个指标之间存在根本性的权衡。
\end{warningbox}

\subsection{推理与训练的核心差异}

\begin{itemize}
    \item \textbf{训练}：可以看到所有 token，可以在序列维度上并行化（矩阵乘法），\textbf{计算受限}（compute-limited）
    \item \textbf{推理}：必须逐 token 顺序生成（每个 token 依赖之前的所有 token），难以充分利用计算资源，\textbf{内存受限}（memory-limited）
\end{itemize}

这一根本差异决定了推理优化的方向：与训练追求更高计算利用率不同，推理优化的核心是减少内存传输。

\begin{figure}[H]
\centering
\begin{tikzpicture}[
    box/.style={draw, rounded corners, minimum width=3.2cm, minimum height=1cm, align=center, font=\small},
    arrow/.style={->, thick, >=stealth}
]
\node[box, fill=blue!15] (train) at (0, 0) {\textbf{训练}\\看到所有 token\\可并行化};
\node[box, fill=red!15] (infer) at (6, 0) {\textbf{推理（生成）}\\逐 token 生成\\顺序执行};
\node[box, fill=green!15] (compute) at (0, -2.2) {\textbf{计算受限}\\GPU 利用率高};
\node[box, fill=orange!15] (memory) at (6, -2.2) {\textbf{内存受限}\\GPU 利用率低};

\draw[arrow] (train) -- (compute);
\draw[arrow] (infer) -- (memory);
\draw[arrow, dashed, gray] (train) -- node[above, font=\footnotesize, gray]{核心差异} (infer);
\end{tikzpicture}
\caption{训练与推理的核心差异：计算受限 vs. 内存受限}
\end{figure}

\subsection{推理优化技术全景}

本节课将覆盖以下推理优化技术，按照从理解问题到解决问题的顺序展开：

\begin{figure}[H]
\centering
\begin{tikzpicture}[
    level/.style={draw, rounded corners, minimum width=4cm, minimum height=0.8cm, align=center, font=\small},
    arrow/.style={->, thick, >=stealth}
]
\node[level, fill=gray!10] (base) at (0, 0) {理解推理工作负载};
\node[level, fill=yellow!15] (lossy1) at (-4, -1.5) {缩减 KV Cache};
\node[level, fill=yellow!15] (lossy2) at (0, -1.5) {替代架构};
\node[level, fill=yellow!15] (lossy3) at (4, -1.5) {量化与剪枝};
\node[level, fill=green!15] (lossless) at (-3, -3) {投机采样（无损）};
\node[level, fill=blue!15] (system) at (3, -3) {系统优化};

\draw[arrow] (base) -- (lossy1);
\draw[arrow] (base) -- (lossy2);
\draw[arrow] (base) -- (lossy3);
\draw[arrow] (base) -- (lossless);
\draw[arrow] (base) -- (system);

\node[font=\footnotesize, gray] at (-4, -0.7) {有损};
\node[font=\footnotesize, gray] at (0, -0.7) {有损};
\node[font=\footnotesize, gray] at (4, -0.7) {有损};
\end{tikzpicture}
\caption{本节课推理优化技术的组织结构}
\end{figure}

\subsection{推理生态}

从事推理优化的公司和项目包括：

\begin{itemize}
    \item \textbf{闭源模型提供商}：OpenAI、Anthropic、Google 等
    \item \textbf{开源权重模型服务商}：Together、Fireworks、DeepInfra 等
    \item \textbf{开源推理框架}：
    \begin{itemize}
        \item vLLM（UC Berkeley）
        \item TensorRT-LLM（NVIDIA）
        \item TGI（Hugging Face）
    \end{itemize}
\end{itemize}

\subsection{本章小结}

推理是大语言模型从训练到实际部署的关键环节。它不仅出现在最终用户交互中，还贯穿于模型评估、测试时计算和强化学习训练的全过程。推理的核心特征是\textbf{顺序生成}和\textbf{内存受限}，这决定了后续所有优化技术的方向。

%% ========== Section 2: 理解推理工作负载 ==========

